\scanpage{15}
% archive foliation: 14
\begin{center}- 99 -\end{center}\hfill n\textsuperscript{o} 268

points génériques indépendants de $A$, on a
\[
s_m(u_1,\ldots,u_m)=u_1+\cdots+u_m\,.
\]
Onnotera $p_i:A_{(m)}\longrightarrow A_i$ la projection sur le i-ème facteur. Si $X$ est un diviseur sur $A$, alors $p_i^{-1}(X)$ est égal au diviseur $A\times\cdots\times X\times\cdots\times A$, avec $X$ au lieu de $A$ au i-ème facteur.

\uline{Théorème} 6 - Soit $X$ un diviseur sur $A$. Alors $X\equiv0$ si et seulement si $s_m^{-1}(X)\sim\sum_{i=1}^m p_i^{-1}(X)$.

\uline{Démonstration} - Pour simplifier les notations, prenons $m=2$. Pour $u$ générique sur $A$, on a
\[
s_2^{-1}(X).(A\times u)=X_{-u}\times u
\]
et par conséquent
\[
[s_2^{-1}(X)-X\times A].(A\times u)=(X_{\unsure{-1|-u}}-X)\times u.
\]
Si $X\equiv0$, on a d'après le principe de la balançoire (Appendix, th. 6)
\[
s_2^{-1}(X)-X\times A\sim A\times Y
\]
avec un diviseur $Y$ convenable. Par symétrie, on doit avoir $Y\sim X$. La réciproque est évidente.

\uline{Corollaire} - Soient $\alpha:A\longrightarrow B$ et $\beta:A\longrightarrow B$ deux homomorphismes d'une variété abélienne $A$ dans une variété abélienne $B$, et soit $X$ un diviseur sur $B$, $X\equiv0$, et soit $\xi=\operatorname{Cl}(X)$. Soient $\alpha^{-1}$, $\beta^{-1}$ les applications induites de $D(B)/D_\ell(B)$ dans $D(A)/D_\ell(A)$. Alors on a
\[
(\alpha+\beta)^{-1}(\xi)=\alpha^{-1}(\xi)+\beta^{-1}(\xi).
\]
En particulier, si $A=B$ et $n$ est un entier, on a
\[
(n\delta)^{-1}(\xi)=n\xi.
\]
\uline{Démonstration} . - Soit $\lambda=(\alpha,\beta)$ l'application composée de $A$ dans $B\times B$, i.e. L'application $\lambda(u)=(\alpha(u),\beta(u))$, et soit $s_2$ la somme sur $B$. On a $\alpha+\beta=s_2\lambda$.
\diagram{\begin{tikzcd}A\arrow[r,"\lambda"]&B\times B\arrow[r,"s_2"]&B\end{tikzcd}}
